\chapter{Abscess (Skin Abscess)}

\section*{Definition}
A skin abscess is a localized collection of pus within the dermis or deeper soft tissue, usually caused by a bacterial infection. It is a walled-off pocket that often requires incision and drainage; some small abscesses may drain spontaneously, and not every inflamed cyst is an infection.

\section*{When to See a Doctor}
\begin{itemize}
\item Tender, red, swollen, or hardened lump under the skin
\item Warmth and redness of the overlying skin
\item Fluctuance (a clinician may feel fluid beneath the skin)
\item Pain, tenderness, or spontaneous drainage of pus
\item Fever, chills, rapid heartbeat, weakness, or feeling very unwell
\item Redness spreading away from the lump or enlarged nearby lymph nodes
\item Urgent assessment is particularly important for abscesses on the face, near the eye, hand, breast, genital or perianal area, or in a person with diabetes or impaired immunity
\end{itemize}

\section*{How It Develops}
Bacteria, most commonly Staphylococcus aureus, can enter through a hair follicle, minor wound, injection site, or other break in the skin. The immune response brings inflammatory cells to the site; bacteria, immune cells, and tissue debris form pus. The surrounding inflammation may localize the infection, while increasing pressure causes pain and fluctuation. Some lesions called cysts are inflamed because their contents have leaked into surrounding tissue and may not require antibiotics.

\section*{Causes}
\begin{itemize}
\item Bacterial infection, most often Staphylococcus aureus, including methicillin-resistant S. aureus (MRSA)
\item Mixed skin or gastrointestinal bacteria in abscesses at certain sites, such as perineal or perirectal areas, or after penetrating injury
\item Blocked or infected hair follicle, minor wound, injection site, or foreign body
\item Hidradenitis suppurativa, which causes recurrent inflammatory nodules and draining tunnels in skin folds and is not simply a series of routine bacterial abscesses
\item Diabetes, impaired circulation, edema, or immunosuppression, which increase the risk of infection or poor healing rather than directly causing every abscess
\end{itemize}

\section*{Related Symptoms}
\begin{itemize}
\item Cellulitis, a spreading infection of the surrounding skin
\item Boil (furuncle) or carbuncle, which involves one or more infected hair follicles
\item Hidradenitis suppurativa, a chronic inflammatory skin disease with recurrent nodules, abscess-like lesions, and draining tunnels
\item Fever or sepsis if infection spreads beyond the local area
\end{itemize}

\section*{Medical Evaluation}
\begin{itemize}
\item Clinical examination assesses the size, depth, location, fluctuance, surrounding cellulitis, pain, neurovascular status, and signs of systemic illness.
\item Point-of-care or formal ultrasound can help distinguish an abscess from cellulitis or another lump, especially when the lesion is deep, small, or uncertain on examination.
\item Pus culture and susceptibility testing are recommended in many drained abscesses, and are particularly useful for recurrent abscesses, treatment failure, severe infection, or impaired immunity. A typical uncomplicated abscess may sometimes be treated without culture according to local practice.
\item Blood cultures are not routine for a simple skin abscess but may be obtained when there are systemic signs of infection or sepsis.
\end{itemize}

\section*{Treatment Options}
\begin{itemize}
\item Incision and drainage is the primary treatment for most fluctuant abscesses and should be performed by a trained clinician with appropriate pain control and sterile technique.
\item Routine packing is not necessary for every small uncomplicated abscess and may increase pain; it may be used selectively for larger, deep, or specialized wounds.
\item Antibiotics are selected by a clinician when there is systemic illness, extensive cellulitis, multiple lesions, impaired host defenses, difficult anatomic location, inadequate drainage, recurrence, or treatment failure. When MRSA coverage is indicated, options may include trimethoprim--sulfamethoxazole or clindamycin according to local resistance, allergies, interactions, and patient factors.
\item Warm compresses and appropriate wound care may help a small superficial lesion or a drained wound.
\item Pain management and follow-up are guided by the size, site, cause, and response to treatment.
\end{itemize}

\section*{Self-Care and Home Management}
\begin{itemize}
\item Apply a clean, warm, moist compress for a small superficial lump while arranging medical advice; do not rely on this if the lesion is enlarging or painful.
\item Do not squeeze, puncture, cut, or attempt to drain the abscess at home.
\item Keep the area clean, cover draining material with a clean dry dressing, wash hands after touching it, and do not share towels or razors.
\item Use over-the-counter pain medicine only as directed and only if it is safe for you.
\item Seek medical care for any suspected abscess, and seek urgent care for fever, spreading redness, severe pain, rapid enlargement, or an abscess in a high-risk location or in a person with diabetes or impaired immunity.
\end{itemize}

\section*{References}
\begin{thebibliography}{99}
\bibitem{abscess:ref1} Stevens DL, Bisno AL, et al. Practice guidelines for the diagnosis and management of skin and soft tissue infections: 2014 update by the Infectious Diseases Society of America. \textit{Clin Infect Dis}. 2014;59(2):e10--e52.
\bibitem{abscess:ref2} MedlinePlus. Skin Abscess. U.S. National Library of Medicine. Reviewed 2024.
\bibitem{abscess:ref3} Centers for Disease Control and Prevention. Preventing and Treating Skin Infections. Updated 2023.
\bibitem{abscess:ref4} Singer AJ, Talan DA. Management of skin abscesses in the era of methicillin-resistant Staphylococcus aureus. \textit{N Engl J Med}. 2014;370(11):1039--1047.
\bibitem{abscess:ref5} Gottlieb M, DeMott JM, Hallock M, et al. Systemic antibiotics for the treatment of skin and soft tissue abscesses: a clinical practice guideline. \textit{BMJ}. 2018;360:k243.
\end{thebibliography}
